\documentclass[tikz,border=6pt]{standalone}
\usepackage{ctex}
\usepackage{amsmath}
\usepackage{pgfplots}
\pgfplotsset{compat=1.18}
\usetikzlibrary{calc,angles,quotes,arrows.meta,positioning,decorations.markings}
\begin{document}
\begin{tikzpicture}[scale=1.9,>=Stealth]
  % 坐标轴
  \draw[->,gray!70] (-0.4,0) -- (2.4,0) node[right,black] {$a$ 轴（$\cos\varphi$ 方向）};
  \draw[->,gray!70] (0,-0.4) -- (0,1.7) node[above,black] {$b$ 轴（$\sin\varphi$ 方向）};

  % 取例题值 a=sqrt3≈1.732, b=1, R=2, φ=30°
  \coordinate (O) at (0,0);
  \coordinate (P) at (1.732,1);

  % 模 R 的圆弧（半径 2）参考
  \draw[dashed,gray!45] (2,0) arc (0:30:2);

  % 向量 (a,b)
  \draw[->,very thick,blue!70!black] (O) -- (P)
        node[midway,above left=-1pt,blue!70!black] {$R=\sqrt{a^2+b^2}$};

  % 分量虚线
  \draw[dashed,blue!40] (P) -- (1.732,0);
  \draw[dashed,blue!40] (P) -- (0,1);

  % a 分量
  \draw[->,red!70!black,thick] (O) -- (1.732,0)
        node[midway,below] {$a$};
  % b 分量
  \draw[->,green!45!black,thick] (1.732,0) -- (P)
        node[midway,right] {$b$};

  % 辐角 φ（从正 a 轴量到向量）
  \path (O) -- (1,0) coordinate (Xdir);
  \draw pic[draw,->,black,"$\varphi$",angle eccentricity=1.5,angle radius=0.62cm]
        {angle=Xdir--O--P};

  % 顶点标注
  \fill[blue!70!black] (P) circle (0.6pt) node[above right=-1pt] {$(a,\,b)$};

  % 半透明注释框
  \node[align=left,fill=yellow!18,fill opacity=0.85,text opacity=1,
        draw=gray!50,rounded corners,inner sep=4pt,font=\small]
        at (1.5,-0.95)
        {$a\sin x+b\cos x=R\sin(x+\varphi)$\\[1pt]
         模 $R=\sqrt{a^2+b^2}$ 即\textbf{振幅}\\[1pt]
         辐角 $\varphi$ 即\textbf{初相}};

  % 本例数值
  \node[font=\footnotesize,gray!55!black] at (2.0,1.32)
        {例：$a=\sqrt3,\,b=1\Rightarrow R=2,\ \varphi=\dfrac{\pi}{6}$};
\end{tikzpicture}
\end{document}
